% MATLAB Practical: Vector manipulation
% This practical follows an introduction to MATLAB coding.
% It is designed for the students to practice:
% - Vectors and vector operations
% - Element-wise calculations
% - Plotting a simple function

\documentclass[12pt]{article}
\usepackage[margin=1in]{geometry}
\usepackage{amsmath}
\usepackage{amssymb}
\usepackage{verbatim}
\usepackage{fancyhdr}

\pagestyle{fancy}
\fancyhf{}
\cfoot{\thepage}

\title{Module 1 Practical 1: MATLAB in Oceanography}
\author{}
\date{Date: 16-02-2026}

\begin{document}
\maketitle

\section*{Instructions}
\begin{itemize}
    \item Complete all questions in a Live Script .
    \item Comment all necessary steps.
    \item Save your script as: \texttt{student\_number }.
    \item Each main question carries 5 marks. The bonus is 2 marks.
\end{itemize}

\section*{Question 1: Wave Velocity Vector Operations (5 marks)}
Given a simple ocean wave with horizontal velocity components $u$ (x-direction) and $v$ (y-direction):

\textbf{Tasks:}
\begin{enumerate}
    \item Define $u$ and $v$ as vectors for a wave with 5 points along the wave crest:
    \[
    u = (0.5, 0.7, 0.6, 0.4, 0.3), \quad
    v = (0.2, 0.3, 0.1, 0.0, -0.1)
    \]
    \item Calculate the \textbf{speed} of the wave at each point using:
    \[
    \text{speed} = \sqrt{u^2 + v^2}
    \]
    \item Display the result within the live editor.
\end{enumerate}

\textbf{Hints:} Use element-wise operations and a function for the square root and display.


\section*{Question 2: Wave Height as a Function of Time (5 marks)}
Consider a simple sinusoidal wave with height $H(t)$ given by:
\[
H(t) = 2 \sin(2 \pi f t)
\]
where $f = 0.5$ Hz, and $t$ is time in seconds.

\textbf{Tasks:}
\begin{enumerate}
    \item Define \texttt{t} as a vector from 0 to 10 seconds in intervals of 0.1 s.
    \item Compute the corresponding wave height vector \texttt{H}.
\end{enumerate}

\textbf{Hints:} Use \verb|sin()| and \verb|pi|. Remember MATLAB operates element-wise on vectors.


\section*{Question 3: Plot Wave Height vs Time (5 marks)}
\textbf{Tasks:}
\begin{enumerate}
    \item Plot wave height \texttt{H} against time \texttt{t}.
    \item Label the axes and add a title.
    \item Set the grid on for clarity.
\end{enumerate}

\section*{Bonus Question (2 marks)}
Retrieve the wave height at a specific time, e.g., \texttt{t = 3.5 s}.

\textbf{Hint:} Use indexing to find the correct element in the vector \texttt{H}.

\begin{verbatim}
Ensure your script is saved correctly and submission is done through Amathuba
\end{verbatim}

\end{document}